\documentclass[tikz,border=8pt]{standalone}
\usepackage{fontspec}
\usepackage{xeCJK}
\IfFontExistsTF{Segoe UI}{\setmainfont{Segoe UI}}{\setmainfont{Arial}}
\IfFontExistsTF{Microsoft JhengHei}{\setCJKmainfont{Microsoft JhengHei}}{\IfFontExistsTF{Noto Sans CJK TC}{\setCJKmainfont{Noto Sans CJK TC}}{\setCJKmainfont{Noto Sans TC}}}
\IfFontExistsTF{Microsoft JhengHei}{\setCJKsansfont{Microsoft JhengHei}}{\IfFontExistsTF{Noto Sans CJK TC}{\setCJKsansfont{Noto Sans CJK TC}}{\setCJKsansfont{Noto Sans TC}}}
\xeCJKsetup{CJKglue={\hskip 0.045em plus 0.015em minus 0.005em}, CJKecglue={\hskip 0.08em plus 0.02em}}
\IfFileExists{../../../FontAwesome.otf}{\newfontfamily\FA[Path=../../../]{FontAwesome.otf}}{\newfontfamily\FA{Segoe UI Symbol}}
\newcommand{\faIcon}[1]{{\FA\symbol{#1}}}
\usetikzlibrary{arrows.meta}
\definecolor{navy}{HTML}{0F172A}
\definecolor{muted}{HTML}{334155}
\definecolor{paper}{HTML}{FFFDF8}
\definecolor{blue}{HTML}{2563EB}
\definecolor{bluebg}{HTML}{DBEAFE}
\definecolor{green}{HTML}{00856A}
\definecolor{greenbg}{HTML}{DDFCEB}
\definecolor{gold}{HTML}{D97706}
\definecolor{goldbg}{HTML}{FEF3C7}
\definecolor{line}{HTML}{CBD5E1}
\begin{document}
\begin{tikzpicture}[x=1cm,y=1cm,>=Latex]
\path[fill=paper] (0,0) rectangle (16,9.6);
\node[font=\bfseries\fontsize{18}{22}\selectfont,text=navy,align=center,text width=14.4cm] at (8,8.82) {母語是橋，不是另一間教室};
\node[font=\fontsize{10.4}{13}\selectfont,text=muted,align=center,text width=13.8cm] at (8,8.12) {支援的方向，是讓學生走回共同討論，而不是留在旁邊};
\tikzset{
  stagebox/.style={rounded corners=8pt,line width=1pt,align=center,text width=2.55cm,minimum height=1.14cm,font=\sffamily\fontsize{8.7}{10.8}\selectfont,text=navy}
}
\node[stagebox,draw=blue,fill=bluebg] (home) at (2.45,5.1) {\bfseries 母語抓概念\\先聽懂};
\node[stagebox,draw=gold,fill=goldbg] (term) at (5.95,5.1) {\bfseries 對照專業詞\\知道邊界};
\node[stagebox,draw=green,fill=greenbg] (shared) at (9.45,5.1) {\bfseries 回到課堂語言\\能參與};
\node[stagebox,draw=blue,fill=bluebg] (contrib) at (12.95,5.1) {\bfseries 修正詞表\\也能貢獻};
\foreach \a/\b in {home/term,term/shared,shared/contrib} {\draw[->,draw=navy,line width=1.1pt] (\a)--(\b);}
\draw[->,draw=green,line width=1pt] (contrib.south) .. controls (10.4,2.45) and (5.65,2.45) .. (home.south);
\node[rounded corners=8pt,draw=line,fill=white,line width=.9pt,align=center,text width=8.8cm,minimum height=.82cm,font=\fontsize{9.2}{11.5}\selectfont,text=muted] at (8,2.22) {學生不是被照顧的對象，也可以共同建造課程語言。};
\node[font=\fontsize{10.2}{12.8}\selectfont,text=navy,align=center,text width=13.6cm] at (8,.9) {真正的平等，不是假裝大家起點相同，而是讓更多人走得到同一張桌子。};
\end{tikzpicture}
\end{document}
